% REDTEAM-STRENGTHEN REPORT::REDTEAM-STRENGTHEN13::2026-10-13T15:17:00+09:00
% !TEX program = xelatex
\documentclass[11pt]{article}
\usepackage[margin=0.88in]{geometry}
\usepackage{amsmath,amssymb,amsthm,booktabs,hyperref,longtable,xurl}
\newtheorem{theorem}{Theorem}
\newtheorem{lemma}[theorem]{Lemma}
\newtheorem{proposition}[theorem]{Proposition}
\newcommand{\Res}{\mathsf{Res}}
\newcommand{\sCP}{\mathsf{sCP}}
\newcommand{\EXACT}{\textbf{PROVED/EXACT}}
\newcommand{\KNOWN}{\textbf{KNOWN THEOREM}}
\newcommand{\COMP}{\textbf{COMPUTATIONAL}}
\newcommand{\COND}{\textbf{CONDITIONAL}}
\newcommand{\CAND}{\textbf{CANDIDATE}}
\newcommand{\OPEN}{\textbf{OPEN}}
\newcommand{\RET}{\textbf{FALSIFIED/RETIRED}}
\hypersetup{colorlinks=true,urlcolor=blue,linkcolor=blue,citecolor=blue}
\title{REDTEAM-STRENGTHEN13:\\Strictifying the Explicit DPLL Certificate Bridge}
\author{Independent Red Team}
\date{13 October 2026}

\begin{document}
\sloppy
\maketitle

\begin{center}
\fbox{\parbox{0.94\linewidth}{
\textbf{Result.}  The newest CERTDPLL-BRIDGE1 claim contains a real but
repairable accounting seam.  Berger--Lawrence--Nordvall Forsberg--
Seisenberger count subsumption leaves at zero cost, so directly emitting one
semantic-Cutting-Planes line for every such leaf does \emph{not} justify the
claimed $n+O(|H|)$ bound.  A redundant-split family falsifies that direct
materialization.  A subclause-preserving strictification removes all free
subsumption and yields at most $n$ genuine Resolution inferences from original
input clauses.  Therefore the restricted bridge from an \emph{explicit,
tree-encoded DPLL derivation with one record per rule occurrence} to a
polynomial-bit semantic-CP certificate is \textbf{PROVED/EXACT}.  Pointer,
macro, or DAG-compressed certificates are outside this theorem.  No arbitrary
one-bit SAT decider is converted to such a derivation.  P versus NP remains
\textbf{OPEN}.}}
\end{center}

\section{Claim card and selected barrier}
\begin{description}
\item[] \textbf{Claim ID:} CERTDPLL-STRICT1.
\item[] \textbf{Audited predecessor:} LSPV-4 report
\path{research/agents/lspv-4/2026-10-13_1455_LSPV-4_CERTDPLL-BRIDGE1.tex}.
\item[] \textbf{Input object:} a CNF $H$ and an explicit tree-encoded DPLL
derivation $\Pi$ in the calculus of Berger et al., with one record for each
rule occurrence.  Let $n$ be the number of charged DPLL rule occurrences and
$B$ the bit length of the complete certificate.
\item[] \textbf{Output object:} a strict binary-Resolution derivation using
only clauses of $H$ as axioms, followed by its clause-inequality embedding in
semantic Cutting Planes.
\item[] \textbf{Exact quantifier:} for every such $(H,\Pi)$, the translator
returns a strict proof whose conclusion is a subclause of the clause promised
by the free-subsumption translation.  In particular, for an unsatisfiability
derivation from the empty valuation, the conclusion is the empty clause.
\item[] \textbf{Not claimed:} for every uniform SAT decider there exists, or
can be extracted, a polynomial-size DPLL, Resolution, or semantic-CP proof.
\end{description}

The two major obstacles were (i) the mismatch between zero-cost subsumption
in the cited Resolution calculus and charged semantic-CP output lines, and
(ii) the encoding ambiguity between expanded rule occurrences and a succinct
pointer/macro representation.  The first is repaired below; the second is an
explicit scope boundary.

\section{Primary-literature and model preflight}
Berger et al. define Resolution with a zero-cost subsumption rule and a
charged binary resolution rule.  Their published Theorem 3.7 states that for
every consistent valuation $\Gamma$, if
$\Gamma\vdash_{\mathrm{DPLL}}^n\Delta$, then
$\Delta\vdash_{\Res}^m\overline\Gamma$ for some $m\le n$
\cite{berger}.  Thus $m$ counts Resolution inferences, not the number of
displayed subsumption leaves.  This distinction is exactly the audited seam.

The independent ordinary-machine endpoint was also rechecked in the primary
source.  Fortnow--Lipton--van Melkebeek--Viglas Corollary 5.8 is
\[
 \operatorname{NTIME}[N]\not\subseteq
 \operatorname{DTISP}\!\left[
 N^c,o\!\left(N^{\frac12((c+2)/c^2-c)}\right)\right],\qquad c>1. \tag{1}
\]
Their model is a random-access Turing machine with indexed read-only input
and write-only one-directional output; input and output tapes are excluded
from workspace \cite{flmv}.  At $c=6/5$, the space exponent is $23/45$.
Equation (1) is not used to prove the certificate bridge.  The retired
unrestricted endpoint $t+s<2\cos(\pi/7)$ is not used.

\section{The free-subsumption counterexample}
Take the one-clause formula $H=\{\emptyset\}$.  Above its conflict leaf,
insert $r$ logically useless consistent splits.  The Berger translation has
\[
       n=2^r-1\quad\hbox{charged Resolution inferences},
       \qquad 2^r\quad\hbox{free subsumption leaves}.      \tag{2}
\]
A direct line-by-line semantic-CP materialization has $2^{r+1}-1$ lines,
whereas the asserted budget $n+|H|$ is $2^r$.  The first counterexample is
$r=1$.  Thus the following interpretation is \RET{}:
\begin{quote}
Every leaf of Berger's free-subsumption proof can be emitted as a separate
semantic-CP line while retaining the exact $n+O(|H|)$ line count.
\end{quote}
This does not refute Theorem 3.7 and does not refute the repaired bridge.

\section{Strictification theorem}
Write clauses as sets of signed literals.  A free-subsumption derivation has
leaves $A\subseteq C$ with $A\in H$, and internal steps
\[
     \frac{C\vee\neg p\qquad D\vee p}{C\vee D}.          \tag{3}
\]

\begin{lemma}[No-weakening subclause invariant]
For every node $v$ of a free-subsumption Resolution derivation with displayed
clause $C_v$, one can construct from original axioms and strict binary
resolution only a clause $A_v\subseteq C_v$.  The construction uses at most
one strict inference for each charged Resolution node below $v$ and has no
greater inference depth.
\end{lemma}

\begin{proof}
At a subsumption leaf $A\subseteq C_v$, return the original axiom $A$; emit
no weakening line.  At an internal node (3), recursively obtain
$A\subseteq C\cup\{\neg p\}$ and $B\subseteq D\cup\{p\}$.
If both $\neg p\in A$ and $p\in B$, emit their genuine resolvent
\[
              (A\setminus\{\neg p\})\cup(B\setminus\{p\})
              \subseteq C\cup D.                         \tag{4}
\]
If $\neg p\notin A$, return $A$ unchanged; then $A\subseteq C\subseteq
C\cup D$.  Symmetrically, if $p\notin B$, return $B$.  Hence an internal
node emits either one strict resolution or no line.  Induction proves the
subclause, count, and depth claims.
\end{proof}

\begin{theorem}[Explicit DPLL certificate to strict Resolution]
Let $\Gamma$ be consistent and let
$\Gamma\vdash_{\mathrm{DPLL}}^n H$ be an explicit tree derivation in the
Berger calculus.  There is a strict Resolution derivation from clauses of
$H$ of a clause $C\subseteq\overline\Gamma$ with at most $n$ derived
resolution nodes.  For $\Gamma=\emptyset$, $C=\emptyset$.
\end{theorem}

\begin{proof}
Apply Berger et al.'s \KNOWN{} translation to obtain a free-subsumption
Resolution derivation of $\overline\Gamma$ with at most $n$ charged
resolution nodes, then apply the lemma.  When $\Gamma=\emptyset$, the only
subclause of $\overline\Gamma=\emptyset$ is $\emptyset$.
\end{proof}

This is a theorem about the supplied certificate, not an invocation of a SAT
oracle and not a search theorem.

\section{Semantic-CP embedding and bit-cost ledger}
Encode a clause $C=P\cup\{\neg x_j:j\in M\}$ as
\[
             \sum_{i\in P}x_i-\sum_{j\in M}x_j
             \ge 1-|M|.                                  \tag{5}
\]
Each original clause of $H$ is emitted at most once.  Every strict Resolution
step is a sound two-parent semantic consequence, so (5) embeds the proof in
semantic CP using at most
\[
                        |H|+n                              \tag{6}
\]
lines.  Variable coefficients lie in $\{-1,0,1\}$ and the right-hand side
has $O(\log(N+2))$ bits, where $N$ is the encoded input length.  Inference
depth is at most $n+O(1)$.

For an implementation, sort literal identifiers and maintain a table from
original clauses to their emitted line indices.  A deliberately conservative
simultaneous bound for reading a $B$-bit explicit certificate, constructing
all subclauses, parent indices and pivots, and writing all output bits is
\[
 O\!\left((B+(|H|+n)N)\log(B+N+2)\right)                 \tag{7}
\]
time and the same order of working bits if all intermediate clauses are
retained.  Streaming can improve space, but no improvement is needed here.
If a uniform proof-emitting machine runs in $T_A(H)$ bit steps and its
write-only output is charged, then $B,n\le T_A(H)$, so (7) is polynomial in
$N+T_A(H)$.  Discovery, rejected branches, repeated rule occurrences,
certificate parsing, replay, buffers, clause dictionaries, parent indexes,
ancestor state, and output are therefore not free.  Time and space refer to
this same translator execution.

The result is \EXACT{} for the explicit tree-certificate class.  It is not a
claim that arbitrary polynomial-time deciders emit certificates.

\section{Encoding firewall: trees are not pointer macros}
The size $n$ in Berger et al. is defined inductively by expanded rule
occurrences.  Consider the macro
\[
        P_r=\operatorname{Split}(y_r,P_{r-1},P_{r-1}).    \tag{8}
\]
A pointer representation has $r+1$ records, whereas its expanded tree has
$2^{r+1}-1$ occurrences.  At $r=30$ these are respectively $31$ and
$2{,}147{,}483{,}647$.  Therefore this report does not silently substitute
pointer-record count for theorem size.  A DAG- or macro-compressed DPLL
certificate requires a separate, model-matched simulation theorem.  The
same warning applies to clause-learning solvers: proof logging, restarts,
learned clauses and deletion require their own proof-system analysis.

\section{Bounded reproducible falsification attempt}
The deterministic program
\path{research/agents/proof/2026-10-13_1517_REDTEAM-STRENGTHEN13-test.py}
ran under a 24-second alarm and completed in $0.676629$ seconds.  It performed
$26{,}244$ exhaustive local subclause-rule checks; strictified $1{,}200$
seeded valid free-subsumption trees containing $66{,}255$ charged free
resolution inferences and $67{,}455$ subsumption leaves; and verified every
reachable strict node.  It also checked redundant-split depths $1$ through
$20$ and pointer-macro depths $1$ through $30$.

The exact result file is
\path{research/agents/proof/2026-10-13_1517_REDTEAM-STRENGTHEN13-results.json}
with semantic digest
\texttt{993d295cad0c280269a27747b4907b9930e3579383299ce3ab6a7adde5527ffe}.
Among the random trees, $21{,}059$ local strict resolutions were created and
$45{,}196$ pivot-missing steps were skipped; later pruning left $5{,}848$
reachable strict nodes.  The finite run is \COMP{} evidence for the
implementation.  The symbolic lemma supplies the universal proof.

\section{Hostile controls and proof barriers}
\begin{longtable}{p{0.25\linewidth}p{0.67\linewidth}}
\toprule
Control & Verdict \\
\midrule
Naive free-leaf emission & \RET{} by (2); repaired by strictification.\\
Pointer/macro certificate & Outside the theorem; exponential expansion in
(8) blocks identification with occurrence count.\\
Arbitrary one-bit decider & \OPEN{}; the decider need not expose a search tree,
conflict clause, or proof log.\\
Uniformity and encoding & \EXACT{} for the explicit translator with bound
(7); no unit-cost large integers or free output.\\
SAME\_REALIZATION & Preserved: time and space are charged to one translator
run on one supplied certificate.\\
Conditional $K$ elimination & Unchanged and still conditional; this local
bridge does not remove it.\\
Relativization & The local syntactic translation relativizes only if the
same explicit proof object is supplied; it does not extract such an object
from oracle answers.\\
Algebrization / Natural Proofs & No circuit-property lower bound or
arithmetization is asserted.  These audits do not close the missing
decider-to-certificate implication.\\
\bottomrule
\end{longtable}

\section{Barrier-transfer and post-quantum ledger}
The Berger DPLL-to-Resolution theorem is a \textbf{directly transferable
theorem}: its proof objects and occurrence measure match the first stage.
The strict Resolution-to-semantic-CP map is an \textbf{EXACT OBJECT/THEOREM
TRANSFER} after eliminating free subsumption.  Proof-logging results for
CDCL are only \textbf{ANALOGY/CANDIDATE} until a precise solver model, log
format, replay cost and simulation theorem are fixed.

VLSI fan-out and memory hierarchy correctly warn that a shared pointer is
not a free expanded tree, but supply no lower bound in this proof model; this
is a useful structural analogy requiring a new transfer theorem.
Control-theoretic observability is a misleading analogy here because the
certificate is already explicit.

Schur-square distinguishers, quadratic-relation ranks, Hilbert functions,
syzygies and collision entropy are \textbf{FAILED TRANSFERS}: the audited
object is a syntactic propositional proof, not a code or evaluation map.
Worst-case/average-case extraction and special-soundness collision ideas are
\textbf{ANALOGY/CANDIDATE} for the missing arbitrary-decider bridge, but no
black-box extraction theorem is known here.  Hash or random-oracle arguments
would be conditional and are inadmissible.  The exposed DPLL structure is
algorithmically exploitable by the translator, which is why the theorem is a
restricted positive result rather than a hardness result.

\section{CORPUS REVIEWED}
The run enumerated the current GitHub tree and read root \path{main.tex},
\path{main4.tex}, \path{main5.tex}, \path{AGENTS.md}, \path{CODEX.md}, the
repository research skill and its red-team/research protocols, every current
task specification, FORMAL-SPINE1, the newest verified MAIN, the newest
verified LSPV-1 through LSPV-4 reports, all current Construction/Recovery/ARC/
Red-Team handoffs bearing on the wrapper, and older reports surfaced by
search for DPLL, Resolution, weakening, semantic CP and certificate transfer.
The causal HEC/VSC-HEC/G-PPC-HEC/CRS-PPC/CCSPC branch, governance notes and
the experiment artifacts cited by active claims were also inspected.

The immediate predecessor LSPV-4 artifact CERTDPLL-BRIDGE1 is substantive,
contains its expected header and verification sentinel, and was readable
from GitHub.  No predecessor write recovery was required.

\section{DUPLICATION / CONTRADICTION AUDIT}
CERTDPLL-BRIDGE1 already imported the Berger theorem and proposed the
restricted semantic-CP bridge.  REDTEAM-STRENGTHEN12 already used the same
pivot-disappears subclause idea inside a different spread-lift proof.  The new
contribution is to apply that invariant to the newest DPLL bridge, exhibit
the smallest redundant-split counterexample to the naive line accounting,
prove the corrected exact bound, and enforce the expanded-tree encoding
guard.

The contradiction is local and resolved: ``subsumption has zero proof size''
in the primary calculus does not mean ``its displayed leaf costs zero output
bits'' in a semantic-CP file.  The strictification avoids emitting the leaf
as a weakened clause and restores (6).  No contradiction was found with the
primary theorem itself.

During this run, the repository head changed to commit
\texttt{3be89e9a83f47f...}, whose commit message says that root
\path{main.tex} was restored but whose actual root blob is the 28-byte body
\texttt{PLACEHOLDER\_FOR\_FULL\_CONTENT} (blob prefix
\texttt{725bf15...}).  This is a concurrent malformed-body incident, not a
scientific result.  Root \path{main.tex} is outside this role's write scope
and was left untouched.  The current substantive \path{main4.tex} and
\path{main5.tex} blobs remained available for authorized reconciliation.

\section{WHAT CHANGED RELATIVE TO ROOT MAIN}
The restricted explicit-certificate arrow can be sharpened from
\CAND{} to \EXACT{} only after adding the no-weakening strictification and
tree-occurrence scope guard.  The load-bearing arbitrary uniform-decider to
certificate arrow remains \OPEN{}.  Because the live root file was replaced
concurrently by a malformed placeholder, no root prose was edited or treated
as authoritative for promotion.

\section{ADOPTED / REJECTED / RETIRED PRIOR COMPONENTS}
\begin{description}
\item[] \textbf{Adopted:} Berger et al.'s DPLL-to-free-subsumption-Resolution
theorem with its exact occurrence measure.
\item[] \textbf{Adopted:} REDTEAM-STRENGTHEN12's subclause/no-weakening
pattern, now proved for arbitrary free-subsumption Resolution trees.
\item[] \textbf{Rejected:} identifying zero-cost subsumption with zero-cost
serialized proof lines.
\item[] \textbf{Rejected:} identifying pointer/macro records with expanded
DPLL rule occurrences.
\item[] \textbf{Retired:} the unrestricted
$t+s<2\cos(\pi/7)$ endpoint; only the exact primary Corollary 5.8 statement
in (1) is recorded.
\end{description}

\section{STALE OR OVERSTATED CLAIMS FOUND}
The sentence ``at most $n+O(|H|)$ semantic-CP lines'' is overstated if it
means direct materialization of every node in Berger's free-subsumption
derivation.  It is correct after the strictification theorem, with the sharper
bound $|H|+n$.  Any statement that omits ``explicit tree-encoded, one record
per rule occurrence'' is also too broad.  These repairs do not promote the
global wrapper.

\section{Proof status and next target}
\[
\boxed{
 \begin{array}{l}
 \text{explicit tree-DPLL certificate}\Rightarrow
 \text{strict Resolution}\Rightarrow\text{semantic CP}:\ \EXACT,\\
 \text{arbitrary uniform SAT decider}\Rightarrow
 \text{such a certificate}:\ \OPEN,\\
 \mathrm{P}\stackrel{?}{\ne}\mathrm{NP}:\ \OPEN.
 \end{array}}
\]
The next bounded target is a precise proof-logging CDCL/DRAT-like machine
class: fix its log syntax, charge restarts, deletions and replay, and locate a
primary simulation theorem before asking whether strict no-weakening and bit
depth survive.  The highest-value global target remains the arbitrary
one-bit-decider-to-certificate implication; no current report proves it.

\begin{thebibliography}{9}
\bibitem{berger}
U. Berger, A. Lawrence, F. Nordvall Forsberg, and M. Seisenberger,
``Extracting verified decision procedures: DPLL and Resolution,''
\emph{Logical Methods in Computer Science} 11(1:6), 2015.
\url{https://doi.org/10.2168/LMCS-11(1:6)2015}.

\bibitem{flmv}
L. Fortnow, R. Lipton, D. van Melkebeek, and A. Viglas,
``Time-Space Lower Bounds for Satisfiability,''
\emph{Journal of the ACM} 52(6):835--865, 2005.
\url{https://doi.org/10.1145/1101821.1101822}.
\end{thebibliography}

\end{document}
% END_OF_VERIFIED_REPORT::REDTEAM-STRENGTHEN::REDTEAM-STRENGTHEN13
